% =============================================================================
% Chapter 1 -- What GRAIL Is
% =============================================================================
%
% Audience prerequisite: none.
%
% Reader takeaway: decide in five minutes whether GRAIL is the right tool
%   for the problem.  Walks away knowing the three usage regimes (standalone,
%   one-shot chain, iterative outer loop) and which one matches their host code.
%
% =============================================================================

\chapter{What GRAIL Is}
\label{ch:what}

\refguide{1.1}{Scientific Context}

GRAIL quantifies one specific error: the error a learned surrogate introduces
at the converged fixed point of a coupled iterative simulation.
That error is not the surrogate's per-call prediction error.
It is that error summed and amplified by every iteration the surrogate
participates in before convergence.
This chapter explains the distinction, identifies the three regimes where
it matters, and gives the reader a decision flowchart.

% =============================================================================
\section{The Coupled-Iteration Problem}
\label{sec:coupled-problem}
% =============================================================================

\refguide{2.2}{The Coupled-Iteration Uncertainty Problem}

Every converged simulation is a fixed point.
In radiative transfer, neutrino transport, chemistry networks, and
implicit time advance, the code iterates radiation against matter
(or species against species) until the two agree.
When one block of that iteration is replaced by a neural emulator,
the emulator's per-call residual $\delta\mathcal{T}$ does not stay localized.
It propagates forward through every subsequent iteration.

The converged-state error is the resolvent sum of all those propagations:
\[
  \delta x^* = (I - D\mathcal{T})^{-1}\,\delta\mathcal{T}(x^*) + O(\|\delta\mathcal{T}\|^2).
\]
The operator $(I - D\mathcal{T})^{-1}$ is the resolvent, and it is
the Green's function of the coupled iteration.
Its norm $A = \|(I - D\mathcal{T})^{-1}\|$ is the amplification factor:
how much a single-call residual is magnified before it reaches the converged state.
Standalone emulator validation measures $\|\delta\mathcal{T}\|$.
It does not measure $\|\delta x^*\|$.
In scattering-dominated regimes, $A$ is large and the gap between the two is large.

A point that is easy to misread the first time.
The resolvent identity is \emph{linear} in $\delta\mathcal{T}$ ---
doubling the per-call residual doubles the predicted equilibrium shift ---
but it sums \emph{all orders} in $D\mathcal{T}$.
Expanding as a geometric series,
$\delta\mathcal{T} + D\mathcal{T}\,\delta\mathcal{T}
 + (D\mathcal{T})^2\,\delta\mathcal{T} + \cdots$,
the sum runs to infinity.
In the regime where GRAIL matters --- cycle spectral radius near 1 ---
the tail terms dominate.
Truncating at the first two terms gives the wrong answer:
for a scalar positive feedback mode $K=0.99$ (or a residual aligned with
that eigenvector), two terms give $1.99\,\delta\mathcal{T}$
while the full sum gives $100\,\delta\mathcal{T}$.
The matrix-free GMRES inside GRAIL sums the full series in one
linear-system call without ever materialising the Jacobian.

% =============================================================================
\section{The Three Regimes}
\label{sec:three-regimes}
% =============================================================================

Three qualitatively different ways to embed a learned surrogate exist.
GRAIL is relevant to exactly one of them.

\medskip

\noindent\textbf{Standalone inference.}
The surrogate produces a prediction once and the prediction is not fed
back into the simulation.
An example is a network that classifies a final image.
The emulator error is the per-call error; no amplification occurs.
GRAIL does not apply.
Standard conformal prediction or held-out validation gives the right
uncertainty estimate.

\medskip

\noindent\textbf{One-shot chain.}
The surrogate sits inside a pipeline where each block runs once and passes
its output forward.
An example is a surrogate stellar-evolution model followed by a one-time
atmosphere computation.
The error propagates through the chain but does not feed back.
GRAIL does not apply.
Chain-rule error propagation or Monte Carlo sampling of the chain is the
right tool.

\medskip

\noindent\textbf{Iterative outer loop.}
The surrogate replaces one block inside an outer loop that converges to
a fixed point.
Every iteration calls the surrogate and passes its output back as the
next iteration's input.
This is the regime of NLTE radiative transfer, CCSN neutrino transport,
stiff chemistry networks, and implicit time advance.
The emulator error is amplified by the resolvent before it reaches the
converged state.
GRAIL applies.

\bigskip

The table below collects the three regimes.

\begin{center}
\small
\begin{tabular}{>{\bfseries}p{0.22\textwidth} p{0.28\textwidth} p{0.35\textwidth}}
\toprule
Regime & Defining property & Correct tool \\
\midrule
Standalone inference
  & Surrogate called once; no feedback
  & Standard conformal prediction; held-out validation \\
One-shot chain
  & Each block runs once; no outer loop
  & Chain-rule propagation; MC sampling \\
Iterative outer loop
  & Surrogate inside a fixed-point loop
  & \textbf{GRAIL} \\
\bottomrule
\end{tabular}
\end{center}

% =============================================================================
\section{A Decision Flowchart}
\label{sec:flowchart}
% =============================================================================

The flowchart below walks through four questions.
Each ``No'' leads to a simpler tool.
All four ``Yes'' answers lead to GRAIL.

\bigskip

\begin{center}
\begin{tikzpicture}[
  decision/.style={diamond, draw, align=center, font=\small,
                   inner sep=3pt, aspect=2.0,
                   minimum width=3.6cm, minimum height=1.0cm},
  action/.style={rectangle, draw, rounded corners=3pt, align=center,
                 font=\small, inner sep=6pt,
                 minimum width=3.6cm},
  grail/.style={rectangle, draw=blue!60!black, fill=blue!8,
                rounded corners=3pt, align=center, font=\small\bfseries,
                inner sep=8pt, minimum width=3.6cm},
  exit/.style={rectangle, draw=gray!70, fill=gray!8, rounded corners=3pt,
               align=center, font=\small, inner sep=6pt,
               minimum width=3.8cm},
  arr/.style={-{Stealth[length=5pt]}, thick},
  lbl/.style={font=\scriptsize},
]

% Nodes
\node[decision] (q1) {Does your simulation\\converge to a\\fixed point?};
\node[decision, below=1.2cm of q1] (q2) {Is one block inside\\that loop a learned\\surrogate?};
\node[decision, below=1.2cm of q2] (q3) {Is the surrogate\\called multiple times\\before convergence?};
\node[decision, below=1.2cm of q3] (q4) {Do you care about\\the \emph{converged-state}\\error, not only\\per-call error?};
\node[grail, below=1.2cm of q4] (grail) {\textbf{GRAIL applies.}\\Use \texttt{resolvent\_solve}.};

% Right-side exits
\node[exit, right=2.4cm of q1] (e1) {Standalone inference.\\Standard conformal\\prediction suffices.};
\node[exit, right=2.4cm of q2] (e2) {No surrogate inside\\the loop.  Standard\\validation suffices.};
\node[exit, right=2.4cm of q3] (e3) {One-shot chain.\\Chain-rule propagation\\or MC sampling.};
\node[exit, right=2.4cm of q4] (e4) {Per-call validation\\already gives the\\error you need.};

% Arrows on main path
\draw[arr] (q1) -- node[lbl,left]{Yes} (q2);
\draw[arr] (q2) -- node[lbl,left]{Yes} (q3);
\draw[arr] (q3) -- node[lbl,left]{Yes} (q4);
\draw[arr] (q4) -- node[lbl,left]{Yes} (grail);

% Arrows to exits
\draw[arr] (q1) -- node[lbl,above]{No} (e1);
\draw[arr] (q2) -- node[lbl,above]{No} (e2);
\draw[arr] (q3) -- node[lbl,above]{No} (e3);
\draw[arr] (q4) -- node[lbl,above]{No} (e4);

\end{tikzpicture}
\end{center}

\bigskip

The NLTE testbed bundled with GRAIL answers all four questions ``Yes'':
the slab iterates radiation and matter to a fixed point;
the APNN replaces the rate-equation solver inside that loop;
the APNN is called once per iteration for tens to hundreds of iterations;
and the science question is the accuracy of the converged populations, not the
accuracy of any single population prediction.

% =============================================================================
\section{What GRAIL Computes}
\label{sec:what-grail-computes}
% =============================================================================

\refguide{2.4}{Resolvent as a Green's Function}

GRAIL evaluates the resolvent $(I - D\mathcal{T})^{-1}$ matrix-free via GMRES.
The user supplies four things.

\begin{itemize}
  \item \texttt{x\_star} --- the emulator's converged fixed point, as a flat
        float64 tensor of shape \texttt{(N,)}.
  \item \texttt{delta\_T} --- the surrogate's per-call residual evaluated at
        \texttt{x\_star}.  Same shape and dtype.
  \item \texttt{host\_jvp} --- a callable that applies the host block's Jacobian
        to a vector.  Opaque; GRAIL never inspects it.
  \item \texttt{emulator\_jvp} --- a callable that applies the emulator block's
        Jacobian to a vector.  Same contract.
\end{itemize}

GRAIL returns the predicted correction $\delta x^*$, the GMRES residual history,
and a condition-number estimate.
Chapter~\ref{ch:statevec} describes the state-vector convention.
Chapter~\ref{ch:wrapping} shows how to construct the two JVP callbacks.
Chapter~\ref{ch:output} describes every field of the result object.

% =============================================================================
\section{What This Manual Covers}
\label{sec:manual-scope}
% =============================================================================

The manual is organized into three parts.

\medskip

\noindent\textbf{Part I: Getting Started} (Chapters 1--2).
Chapter~1 is this chapter.
Chapter~\ref{ch:install} covers installation, the conda environment, the
libomp workaround, and a first green test run.

\medskip

\noindent\textbf{Part II: Using GRAIL} (Chapters 3--6).
Chapter~\ref{ch:statevec} describes the state-vector seam: the
level-major packing convention that every JVP callback must follow.
Chapter~\ref{ch:wrapping} shows how to construct \texttt{host\_jvp} and
\texttt{emulator\_jvp} for a real host code, including the four standard
pitfalls.
Chapter~\ref{ch:output} walks through every field of \texttt{ResolventResult}
and shows how to read GMRES convergence diagnostics.
Chapter~\ref{ch:coverage} describes the conformal calibration path: how to
build a residual ensemble, fit a PCA basis, and produce a coverage band with
a distribution-free guarantee.

\medskip

\noindent\textbf{Part III: Going Further} (Chapters 7--10).
Chapter~\ref{ch:debug} provides a symptom-to-cause-to-fix table for the
most common failure modes.
Chapter~\ref{ch:perf} explains matrix-free cost, restart storage,
preconditioning, and reproducible performance measurements.
Chapter~\ref{ch:recipes} runs the supplied radiative-transfer,
electronic-structure, climate and calibration examples.
Chapter~\ref{ch:scf} walks through a second worked application---SCF
electronic structure on stretched H$_2$ with a neural exchange--correlation
surrogate---using the same callback interface, unchanged, against a
structurally distinct host.

\medskip

The GRAIL Reference Guide is the companion document for theory and derivations.
Where this manual needs the math, it points to the Reference Guide by section
rather than reproducing the derivation.
A reader who wants only the ``how'' can work through this manual without the
Reference Guide.
A reader who wants the ``why'' finds it there.
